\newpage

\section{vector 类}
\label{sec:class_vector}

\tkzname{vector} 对象表示从尾点 $A$ 到头点 $B$ 的有向线段。
在内部，它由两个点 \tkzAttr{vector}{za} 和 \tkzAttr{vector}{zb} 构建而成，并将其解释为复数。
该向量同时存储其几何端点和解析数据（分量、长度、方向）。



变量 \tkzVar{vector}{V} 保存一个用于存储向量的表。它是可选的，你可以自由选择变量名。不过，为清晰和一致起见，推荐约定使用 \code{V}。如果你使用自定义变量（例如 Vectors），则必须手动初始化它。函数 \code{init\_elements()} 在被调用时会重新初始化 \code{V} 表。

\vspace{1em}
事实上，与其说它们是严格数学意义上的向量，不如说它们是一类有向线段。

一个向量通过给出两个点（即两个附标）来定义。
|V.AB = vector(z.A, z.B)| 创建向量 $(\overrightarrow{AB})$，即以 $A$ 为起点、表示一个向量的有向线段。这里定义了一些基本运算，例如加法、减法以及数乘。


\subsection{创建向量}
\label{sub:creating_a_vector}

\tkzClass{vector} 类表示两个点之间的向量，或一个给定坐标的自由向量。

\medskip
结果通常存储在 \tkzVar{vector}{V} 中。

\medskip
\textbf{简写：}
\begin{mybox}
\begin{verbatim}
V.v1 = vector:new(z.A, z.B) -- from A to B
简写形式与之等价：
V.v1 = vector(z.A, z.B)
\end{verbatim}
\end{mybox}



\subsection{向量的属性}

\vspace{1em}

  \bgroup
  \small
  \catcode`_=12
  \captionof{table}{向量的属性。}\label{vector:attributes}
  \begin{tabular}{ll}
  \toprule
  \texttt{属性}      & \texttt{参考}\\
  \midrule
  \tkzAttr{vector}{tail}     &  \\
  \tkzAttr{vector}{head}     &  [\ref{ssub:attributes_vector_head}] \\
  \tkzAttr{vector}{type}     &  [\ref{ssub:attributes_vector_type}] \\
  \tkzAttr{vector}{slope}    &  [\ref{ssub:attributes_vector_slope}] \\
  \tkzAttr{vector}{z}        &  [\ref{ssub:attributes_vector_z}] \\
  \tkzAttr{vector}{dx}       &  [\ref{ssub:attributes_vector_d}] \\
  \tkzAttr{vector}{dy}       &  [\ref{ssub:attributes_vector_d}]\\
  \tkzAttr{vector}{norm}     &  [\ref{ssub:attributes_vector_norm}]\\
  \tkzAttr{vector}{mtx}      &  [\ref{ssub:attributes_vector_mtx}] \\
  \bottomrule
  \end{tabular}
  \egroup


\subsubsection{属性 \tkzAttr{vector}{head}}
\label{ssub:attributes_vector_head}

\begin{minipage}{.5\textwidth}
\directlua{
 init_elements()
 z.O = point(0, 0)
 z.A = point(0, 1)
 z.B = point(3, 4)
 z.C = point(1, 2)
 z.D = point(2, 1)
 V.u = vector(z.A, z.B)
 V.v = vector(z.C, z.D)
 V.w = V.u + V.v
 z.E = V.w.head}
\begin{tikzpicture}[gridded]
  \tkzGetNodes
  \tkzLabelPoints(A,B,C,D,O,E)
  \tkzDrawSegments[->,red](A,B)
  \tkzDrawSegments[->,cyan](A,E)
  \tkzDrawSegments[->,blue](C,D B,E)
  \tkzLabelSegment(A,B){$\overrightarrow{u}$}
  \tkzLabelSegment(C,D){$\overrightarrow{v}$}
  \tkzLabelSegment(A,E){$\overrightarrow{w}$}
\end{tikzpicture}
\end{minipage}
\begin{minipage}{.5\textwidth}
\begin{tkzexample}[code only]
\directlua{
 init_elements()
 z.O = point(0, 0)
 z.A = point(0, 1)
 z.B = point(3, 4)
 z.C = point(1, 2)
 z.D = point(2, 1)
 V.u = vector(z.A, z.B)
 V.v = vector(z.C, z.D)
 V.w = V.u + V.v
 z.E = V.w.head}
\end{tkzexample}
\end{minipage}


\subsubsection{属性 \tkzAttr{vector}{type}}
\label{ssub:attributes_vector_type}
沿用前面的数据：

\begin{mybox}
\code{|V.u.type = 'vector'|}
\end{mybox}

\subsubsection{属性 \tkzAttr{vector}{slope}}
\label{ssub:attributes_vector_slope}
该属性给出支撑某个向量代表的直线的斜率。

\begin{tkzexample}[latex = .5\textwidth]
\directlua{
   z.A = point(-1, 0)
   z.B = point (3, 3)
   V.u = vector(z.A, z.B)
   local d = math.deg(V.u.slope)
   tex.print(utils.format_number(d,3))}
\end{tkzexample}

\subsubsection{属性 \tkzAttr{vector}{dx}、\tkzAttr{vector}{dy}}
\label{ssub:attributes_vector_d}

这些属性给出向量在参考基下的坐标。

\begin{tkzexample}[latex=.5\textwidth]
\directlua{
  z.A = point(0, 1)
  z.B = point(3, 4)
  z.C = point(3, 2)
  z.D = point(4, 1)
  V.u = vector(z.A, z.B)
  V.v = vector(z.C, z.D)
  V.w = V.u + V.v
  z.E = V.w.head
  local pc = string.char(37)
  local format = "("..pc.."g,"..pc.."g)"
  tex.print("(dx,dy) =
     ",string.format(format, V.w.dx, V.w.dy))
  }
\begin{center}
  \begin{tikzpicture}[gridded]
  \tkzGetNodes
  \tkzLabelPoints(A,B,C,D,E)
  \tkzDrawSegments[->,red](A,B C,D)
  \tkzDrawSegments[->,blue](A,E B,E)
  \tkzDrawPoints(A,C)
\end{tikzpicture}
  \end{center}
\end{tkzexample}


\subsubsection{属性 \tkzAttr{vector}{norm}}
\label{ssub:attributes_vector_norm}
该属性给出两端点之间线段的长度。

\begin{tkzexample}[latex = .5\textwidth]
\directlua{
   z.A = point(-1, 0)
   z.B = point(3, 3)
   V.u = vector(z.A, z.B)
   V.d = V.u.norm
   tex.print("norm = ",V.d)}
\end{tkzexample}

\subsubsection{属性 \tkzAttr{vector}{mtx}}
\label{ssub:attributes_vector_mtx}
这里将一个以 \code{column vector}（列向量）形式表示的矩阵与所讨论的向量关联起来。

\begin{tkzexample}[latex = .5\textwidth]
\directlua{
   z.O = point(1, 0)
   z.I = point(2, 1)
   V.u = vector(z.O, z.I)
   V.u.mtx:print()
   V.v = V.u.mtx:transpose()
   V.v:print()}
\end{tkzexample}

\subsubsection{属性 \tkzAttr{vector}{z}}
\label{ssub:attributes_vector_z}
该属性对于使用 \code{\^{}} 和 \code{..} 元方法非常有用。

\begin{tkzexample}[latex = .5\textwidth]
\directlua{
 z.A = point(1, 1)
 z.B = point(2, -1)
 z.C = point(0, 1)
 z.D = point(3, 2)
 V.u = vector(z.A, z.B)
 V.v = vector(z.C, z.D)
 n = V.u.z ^ V.v.z
 m = V.u.z .. V.v.z
 tex.print("determinant(u,v) = "..tostring(n))
 tex.print('\\\\')
 tex.print("dot product(u,v) = "..tostring(m))
}
\end{tkzexample}


\newpage

\subsection{vector 类的元方法概览}

\vspace{1em}
\bgroup
\catcode`_=12
\small
\captionof{table}{vector 类的方法。}\label{vector:metamethods}
\begin{tabular}{ll}
\toprule
\texttt{元方法} & \texttt{参考} \\
\midrule
\tkzMeta{vector}{add(u,v)}   & [\ref{ssub:method_vector_add}]  \\
\tkzMeta{vector}{sub(u,v)}   & [\ref{ssub:method_vector_sub}]  \\
\tkzMeta{vector}{unm(u)}     & [\ref{ssub:method_vector_unm}] \\
\tkzMeta{vector}{mul(k,u)}   & [\ref{ssub:method_vector_mul}]        \\
\tkzMeta{vector}{pow(k,u)}   & [\ref{ssub:method_vector_det}]        \\
\tkzMeta{vector}{concat(k,u)}& [\ref{ssub:method_vector_dot}]        \\
\bottomrule
  \end{tabular}
\egroup



\subsection{元方法示例}

\subsubsection{方法 \tkzMeth{vector}{add}}
\label{ssub:method_vector_add}

加法定义如下：

设 V.AB + V.CD 得到一个向量 V.AE，其定义如下

如果 $\overrightarrow{CD} = \overrightarrow{BE} $，则 $\overrightarrow{AB} + \overrightarrow{CD} = \overrightarrow{AB} + \overrightarrow{BE} =\overrightarrow{AE}$

\begin{verbatim}
z.A  = point(0, 1)
z.B  = point(3, 4)
z.C  = point(3, 2)
z.D  = point(4, 1)
V.AB = vector(z.A, z.B)
V.CD = vector(z.C, z.D)
V.AE = V.AB + V.CD  % possible V.AB:add(V.CD)
z.E  = V.AE.head % we recover the final point (head)
\end{verbatim}


\begin{tkzexample}[latex=.5\textwidth]
\directlua{
 init_elements()
 z.A = point(0, 1)
 z.B = point(3, 4)
 z.C = point(3, 2)
 z.D = point(4, 1)
 V.u = vector(z.A, z.B)
 V.v = vector(z.C, z.D)
 V.w = V.u + V.v % or V.u:add(V.v)
 z.E = V.w.head}
\begin{tikzpicture}[gridded]
  \tkzGetNodes
  \tkzLabelPoints(A,B,C,D,E)
  \tkzDrawSegments[->,red](A,B)
  \tkzDrawSegments[->,cyan](A,E)
  \tkzDrawSegments[->,blue](C,D B,E)
  \tkzLabelSegment(A,B){$\overrightarrow{u}$}
  \tkzLabelSegment(C,D){$\overrightarrow{v}$}
  \tkzLabelSegment(A,E){$\overrightarrow{w}$}
\end{tikzpicture}
\end{tkzexample}


\subsubsection{方法 \tkzMeth{vector}{sub}}
\label{ssub:method_vector_sub}

\begin{tkzexample}[latex=.5\textwidth]
\directlua{
 z.A = point(0, 1)
 z.B = point(3, 4)
 z.C = point(3, 2)
 z.D = point(4, 1)
 V.u = vector(z.A, z.B)
 V.v = vector(z.C, z.D)
 V.w = V.u - V.v
 z.E = V.w.head}

\begin{tikzpicture}[gridded]
  \tkzGetNodes
  \tkzLabelPoints(A,B,C,D,E)
  \tkzDrawSegments[->,red](A,B)
    \tkzDrawSegments[->,cyan](A,E)
   \tkzDrawSegments[->,blue](C,D B,E)
  \tkzLabelSegment(A,B){$\overrightarrow{u}$}
  \tkzLabelSegment(C,D){$\overrightarrow{v}$}
  \tkzLabelSegment(A,E){$\overrightarrow{w}$}
\end{tikzpicture}
\end{tkzexample}


\subsubsection{方法 \tkzMeth{vector}{mul}}
\label{ssub:method_vector_mul}
这当然就是数乘。

\begin{tkzexample}[latex=.5\textwidth]
\directlua{
 init_elements()
 z.A = point(0, 1)
 z.B = point(3, 4)
 z.C = point(3, 2)
 z.D = point(4, 1)
 V.u = vector(z.A, z.B)
 V.v = vector(z.C, z.D)
 V.w = V.u + 2 * V.v
 z.E = V.w.head}

\begin{tikzpicture}[gridded]
  \tkzGetNodes
  \tkzLabelPoints(A,B,C,D,E)
  \tkzDrawSegments[->,red](A,B)
    \tkzDrawSegments[->,cyan](A,E)
   \tkzDrawSegments[->,blue](C,D B,E)
  \tkzLabelSegment(A,B){$\overrightarrow{u}$}
  \tkzLabelSegment(C,D){$\overrightarrow{v}$}
  \tkzLabelSegment(A,E){$\overrightarrow{w}$}
\end{tikzpicture}
\end{tkzexample}


\subsubsection{方法 \tkzMeth{vector}{unm}}
\label{ssub:method_vector_unm}

这个元方法让你可以书写 \code{V.w = -V.v}
\begin{tkzexample}[latex=.5\textwidth]
\directlua{
  z.A = point(0, 1)
  z.B = point(3, 3)
  z.C = point(3, 2)
  z.D = point(4, 1)
  V.u = vector(z.A, z.B)
  V.v = vector(z.C, z.D)
  V.v = -V.v
  V.w = V.u + V.v
  z.E = V.w.head}
\begin{center}
  \begin{tikzpicture}[gridded]
  \tkzGetNodes
  \tkzLabelPoints(A,B,C,D,E)
  \tkzDrawSegments[->,red](A,B)
  \tkzDrawSegments[->,cyan](A,E)
  \tkzDrawSegments[->,blue](D,C B,E)
\end{tikzpicture}
\end{center}
\end{tkzexample}

\subsubsection{方法 \tkzMeth{vector}{\textasciicircum{}}}
\label{ssub:method_vector_det}
这里不是指幂运算（那在这里没有意义），而是指两个向量的行列式。注意：所使用的参考系是正交且右手的。

\begin{mybox}
|V.u = vector(z.A, z.B) with z.A = point(xa, ya) and z.B = point(xb, yb)|
|V.v = vector(z.C, z.D) with z.C = point(xc, yc) and z.D = point(xd, yd)|
then |V.u ^ V.v = xa * yb -  xb * ya|
\end{mybox}

注：|u ^ v = u.norm * v.norm * sin(u,v)|

\subsubsection{方法 \tkzMeth{vector}{..}}
\label{ssub:method_vector_dot}
这里不是指拼接（那在这里没有意义），而是指两个向量的点积。注意：所使用的参考系是正交且右手的。

\begin{mybox}
|V.u = vector(z.A, z.B) with z.A = point(xa, ya) and z.B = point(xb, yb)|
|V.v = vector(z.C, z.D) with z.C = point(xc, yc) and z.D = point(xd, yd)|
then |V.u .. V.v = xa * xb +  ya * yb|
\end{mybox}

注：|V.u .. V.v = V.u.norm * V.v.norm * cos(V.u,V.v)|


我们将使用最后这两个方法。利用在本宏包所使用的右手正交参考系中表示出的点积和行列式，我们可以求出该角的余弦和正弦。

\begin{minipage}{.5\textwidth}
\directlua{
 z.O = point(0, 0)
 z.A = point(5, 0)
 z.B = point(1, -4)
 V.u = vector(z.O, z.A)
 V.v = vector(z.O, z.B)
 local dp = V.u .. V.v % dot product
 local d = V.u ^ V.v % determinant
 % costheta = dp / (u.norm * v.norm)
 % sintheta = d / (u.norm * v.norm)
  an = math.atan(d / dp)}
  \begin{tikzpicture}
  \tkzGetNodes
  \tkzDrawSegments[>=stealth,->](O,A O,B)
  \tkzDrawPoints(O)
  \tkzLabelPoints(A,B)
  \tkzLabelPoints[left](O)
  \tkzMarkAngle[<-,>=stealth](B,O,A)
  \tkzLabelAngle[pos=2](B,O,A){$\directlua{
  local format = string.char(37) .. ".2f"
  tex.print(string.format( format , math.deg(an)))}$ deg}
  \end{tikzpicture}
\end{minipage}
\begin{minipage}{.5\textwidth}
\begin{tkzexample}[code only]
\directlua{
  z.O = point(0, 0)
  z.A = point(5, 0)
  z.B = point(1, -4)
  V.u = vector(z.O, z.A)
  V.v = vector(z.O, z.B)
  local dp = V.u .. V.v % dot product
  local d = V.u ^ V.v % determinant
  % costheta = dp / (u.norm * v.norm)
  % sintheta = d / (u.norm * v.norm)
  an = math.atan(d / dp)}
\end{tkzexample}
\end{minipage}

显示角度测量值所需的代码如下：

\begin{tkzexample}[code only]
  \tkzLabelAngle[pos=2](B,O,A){$\directlua{
  local format = string.char(37) .. ".2f"
  tex.print(string.format( format , math.deg(an)))}$ deg}
\end{tkzexample}

% subsubsection method_vector_dot（结束）`



\smallskip
  \tkzClass{vector} 类所提供的主要方法总结如下。

\medskip
\small
\captionof{table}{vector 类的方法。}\label{vector:methods}
\begin{tabular}{ll}
\toprule
\texttt{方法} & \texttt{参考} \\
\midrule
\texttt{构造函数} & \\
\midrule
\tkzMeth{vector}{new(pt, pt)} & \ref{sub:creating_a_vector} \\
\midrule
\texttt{返回布尔值的方法} & \\
\midrule
\tkzMeth{vector}{is\_zero([EPS])} & \\
\tkzMeth{vector}{is\_parallel(v,[EPS])} & \\
\tkzMeth{vector}{is\_orthogonal(v,[EPS])} & \\
\midrule
\texttt{返回点的方法} & \\
\midrule
\tkzMeth{vector}{get()}   & [\ref{sub:vector_other_methods}] \\
\midrule
\texttt{返回数值的方法} & \\
\midrule
\tkzMeth{vector}{dot(v)}        & [\ref{sub:vector_other_methods}] \\
\tkzMeth{vector}{cross(v)}      & [\ref{sub:vector_other_methods}] \\
\tkzMeth{vector}{angle\_to(v)}  & [\ref{sub:vector_other_methods}] \\
\midrule
\texttt{返回向量的方法} & \\
\midrule
\tkzMeth{vector}{add(v)}        & [\ref{sub:vector_other_methods}] \\
\tkzMeth{vector}{scale(d)}      & [\ref{sub:vector_other_methods}] \\
\tkzMeth{vector}{normalize()}   & [\ref{ssub:method_vector_normalize}]     \\
\tkzMeth{vector}{at (pt)}       & [\ref{ssub:method_vector_at}]            \\
\tkzMeth{vector}{rotate($\theta$)} & [\ref{sub:vector_other_methods}] \\
\tkzMeth{vector}{orthogonal([side],[length])} & [\ref{ssub:method_vector_orthogonal_r}] \\
\bottomrule
\end{tabular}



\subsection{返回布尔值}

\subsubsection{判定谓词：\tkzMeth{vector}{is\_zero}、\tkzMeth{vector}{is\_parallel}、\tkzMeth{vector}{is\_orthogonal}}

\begin{verbatim}
if V:is_zero() then ... end
if V:is_parallel(W) then ... end
if V:is_orthogonal(W) then ... end
\end{verbatim}

这些方法检验向量之间的基本几何关系，使用一个可选的
容差 \tkzKey{EPS}，默认为 \tkzSysVar{tkz.epsilon}。它们很方便，
尤其在数值计算中需要判定两个方向是否应视为平行或正交时。

\subsection{返回向量}


\subsubsection{方法 \tkzMeth{vector}{normalize()}}
\label{ssub:method_vector_normalize}
该方法产生一个与原向量共线的归一化（单位）向量。

 \begin{tkzexample}[latex = .5\textwidth]
   \directlua{
   init_elements()
   z.A = point(0, 1)
   z.B = point(3, 4)
   V.AB = vector(z.A, z.B)
   V.AN = V.AB:normalize()
   z.N = V.AN.head}
   \begin{tikzpicture}
   \tkzGetNodes
   \tkzDrawSegments(A,B)
   \tkzDrawPoints(A,B,N)
   \tkzLabelPoints(A,B,N)
   \end{tikzpicture}
 \end{tkzexample}

\subsubsection{方法 \tkzMeth{vector}{at()}}
\label{ssub:method_vector_at}

 \begin{tkzexample}[latex = .5\textwidth]
   \directlua{
   init_elements()
   z.A = point(0, 1)
   z.B = point(3, 4)
   z.O = point(0, 0)
   V.AB = vector(z.A, z.B)
   V.OC = V.AB:at(z.O)
   z.C = V.OC.head}
   \begin{tikzpicture}
   \tkzGetNodes
   \tkzDrawSegments(A,B O,C)
   \tkzDrawPoints(A,B,O,C)
   \tkzLabelPoints(A,B,O,C)
   \end{tikzpicture}
 \end{tkzexample}

\subsubsection{方法  \tkzMeth{vector}{orthogonal([side],[length])}}
\label{ssub:method_vector_orthogonal_r}

\begin{verbatim}
Vperp  = V:orthogonal()           -- default: "ccw"
VperpC = V:orthogonal("cw")       -- clockwise orthogonal
Vlen   = V:orthogonal("ccw", 2.0) -- orthogonal of prescribed length
\end{verbatim}

\tkzMeth{vector}{orthogonal} 方法构造一个与 \code{V} 正交、
且尾点相同的向量：
\begin{itemize}
  \item 可选参数 \tkzKey{side} 可取 \code{"ccw"}（逆时针，
        默认值）或 \code{"cw"}（顺时针），
  \item 可选参数 \tkzKey{length} 可用于指定所得向量的模；
        若省略，则长度由当前向量旋转决定。
\end{itemize}

\subsection{其他方法}
\label{sub:vector_other_methods}

\begin{itemize}
  \item \tkzMeth{vector}{get()} 返回向量的尾点和头点。
  \item \tkzMeth{vector}{add(v)} 返回两个向量之和，以
        当前向量的尾点为起点（亦可写作 \code{V + W}）。
  \item \tkzMeth{vector}{scale(d)} 返回乘以 \code{d} 后的向量，
        尾点相同（亦可写作 \code{d * V}）。
  \item \tkzMeth{vector}{dot(v)} 返回点积（亦可写作 \code{V .. W}）。
  \item \tkzMeth{vector}{cross(v)} 返回叉积
        $x_V y_W - y_V x_W$（亦可写作 \code{V \textasciicircum{} W}）。
  \item \tkzMeth{vector}{angle\_to(v)} 返回从当前向量到
        \code{v} 的有向角，取值范围为 $[0, 2\pi)$。
  \item \tkzMeth{vector}{rotate(theta)} 返回绕其尾点旋转
        \code{theta}（弧度）后的向量。
\end{itemize}

\begin{verbatim}
\directlua{
 init_elements()
 z.A = point(0, 0)
 z.B = point(3, 1)
 z.C = point(1, 2)
 V.u = vector(z.A, z.B)
 V.v = vector(z.A, z.C)
 z.S = V.u:add(V.v).head
 z.R = V.u:rotate(math.pi / 2).head
 z.D = V.u:scale(-0.5).head
 tex.print("dot = " .. V.u:dot(V.v) .. ", cross = " .. V.u:cross(V.v)..", angle = " .. tkz.round(math.deg(V.u:angle_to(V.v)), 2) .. " degrees")}
\begin{tikzpicture}
\tkzGetNodes
\tkzDrawSegments[->](A,B A,C A,S A,R A,D)
\tkzDrawPoints(A,B,C,S,R,D)
\tkzLabelPoints(A,B,C,S,R,D)
\end{tikzpicture}
\end{verbatim}

\directlua{
 init_elements()
 z.A = point(0, 0)
 z.B = point(3, 1)
 z.C = point(1, 2)
 V.u = vector(z.A, z.B)
 V.v = vector(z.A, z.C)
 z.S = V.u:add(V.v).head
 z.R = V.u:rotate(math.pi / 2).head
 z.D = V.u:scale(-0.5).head
 tex.print("dot = " .. V.u:dot(V.v) .. ", cross = " .. V.u:cross(V.v)..", angle = " .. tkz.round(math.deg(V.u:angle_to(V.v)), 2) .. " degrees")
   tex.print('\\\\')}
\begin{tikzpicture}
\tkzGetNodes
\tkzDrawSegments[->](A,B A,C A,S A,R A,D)
\tkzDrawPoints(A,B,C,S,R,D)
\tkzLabelPoints(A,B,C,S,R,D)
\end{tikzpicture}

\begin{tkzexample}[latex=.5\textwidth]
 \directlua{
 init_elements()
 z.A = point(0, 1)
 z.B = point(3, 4)
 V.AB = vector(z.A, z.B)
 V.AR = V.AB:orthogonal(2 * math.sqrt(2))
 z.R = V.AR.head
 V.AS = V.AB:orthogonal("cw",2)
 z.S = V.AS.head}
\begin{center}
 \begin{tikzpicture}[gridded]
  \tkzGetNodes
  \tkzDrawSegments[>=stealth,->,
    red](A,B A,R A,S)
  \tkzLabelPoints(A,B,R,S)
\end{tikzpicture}
\end{center}
\end{tkzexample}



\endinput
